\SourceLeaf{85}
\fancyhead[L]{\small Die erdmagnetische Kraft in absoluten Einheiten}
\renewcommand{\Article}[1]{\par\medskip\Needspace{4\baselineskip}\pdfbookmark[1]{Art. #1}{magnetism-art-#1}\centerline{\textbf{#1.}}\nobreak\smallskip}
\pdfbookmark[0]{Die Intensität der erdmagnetischen Kraft in absoluten Einheiten}{magnetism-body}
\par\medskip
\centerline{DIE INTENSITÄT}
\centerline{DER ERDMAGNETISCHEN KRAFT}
\centerline{IN ABSOLUTEN EINHEITEN.}
\medskip
Zur vollständigen Bestimmung der erdmagnetischen Kraft an einem gegebenen Ort sind drei Elemente erforderlich. Das erste ist die Deklination, also der Winkel zwischen der Ebene, in der die Kraft wirkt, und der Meridianebene. Das zweite ist die Neigung ihrer Richtung zur Horizontalebene. Das dritte ist die Intensität. Die Deklination ist bei allen Anwendungen in der Schifffahrt und Geodäsie als das wichtigste Element anzusehen. Sie hat die Beobachter und Physiker von Anfang an beschäftigt, und seit einem Jahrhundert widmen sie auch der Inklination unablässig ihre Aufmerksamkeit. Dagegen blieb das dritte Element, die Intensität, bis in die jüngere Zeit völlig unbeachtet, obwohl es wissenschaftlich ebenso sehr der Untersuchung wert ist. Dem bedeutenden \textsc{Humboldt} gebührt neben seinen vielen anderen Verdiensten auch das Verdienst, dass er als einer der Ersten diesem Gegenstand seine Aufmerksamkeit zuwandte. Auf seinen Reisen sammelte er zahlreiche Beobachtungen zur relativen Intensität des Erdmagnetismus, aus denen hervorging, dass diese Intensität auf dem Weg vom magnetischen Äquator zum Pol stetig zunimmt. Sehr viele Physiker sind den Spuren dieses Naturforschers gefolgt und haben bereits so viele Bestimmungen zusammengetragen, dass der hervorragende und um den Erdmagnetismus hochverdiente \textsc{Hansteen} vor Kurzem den Entwurf einer Weltkarte mit isodynamischen Linien veröffentlichen konnte.

Ihr Verfahren besteht darin, die Zeit zu beobachten, in der dieselbe Magnetnadel an verschiedenen Orten dieselbe Anzahl von Schwingungen ausführt, oder die Anzahl der Schwingungen im selben Zeitintervall zu beobachten. Dabei wird die Intensität dem Quadrat der Anzahl der Schwingungen in einer gegebenen Zeit proportional gesetzt. So werden die Gesamtintensitäten miteinander verglichen, wenn eine in ihrem Schwerpunkt aufgehängte Inklinationsnadel um eine zum magnetischen Meridian senkrechte horizontale Achse schwingt. Die Intensitäten der horizontalen Kraft werden verglichen, wenn eine horizontale Nadel um eine \SourceLeaf{86}vertikale Achse schwingt. Die zweite Beobachtungsweise ist genauer, und aus ihren Ergebnissen lassen sich bei bekannten Inklinationen leicht die Gesamtintensitäten ermitteln.

Das Verfahren beruht offensichtlich auf der Annahme, dass die Verteilung des freien Magnetismus in den Teilchen der für einen solchen Vergleich verwendeten Nadel bei den einzelnen Versuchen unverändert geblieben ist. Hätte nämlich die Magnetkraft der Nadel im Laufe der Zeit etwas nachgelassen, so würde die Nadel gerade deshalb später langsamer schwingen. Ein Beobachter, der von dieser Änderung nichts weiß, würde der Intensität des Erdmagnetismus am späteren Ort einen zu kleinen Wert zuschreiben. Erstrecken sich die Versuche über einen mäßig langen Zeitraum und verwendet man eine sorgfältig magnetisierte Nadel aus gut gehärtetem Stahl, so ist eine erhebliche Schwächung kaum zu befürchten. Die Unsicherheit verringert sich zudem, wenn mehrere Nadeln für die Vergleiche verwendet werden. Schließlich gewinnt die Annahme an Glaubwürdigkeit, wenn sich nach der Reise am ersten Ort zeigt, dass sich die Schwingungsdauer der Nadel nicht geändert hat. Doch welche Vorkehrungen man auch trifft, eine allmähliche Schwächung der Magnetkraft der Nadel ist kaum zu vermeiden. Eine solche Übereinstimmung ist daher nach längerer Abwesenheit selten zu erwarten. Beim Vergleich der Intensitäten für weit voneinander entfernte Orte der Erde wird man deshalb meist nicht die Genauigkeit und Sicherheit erreichen, die man anstreben muss.

Dieser Nachteil fällt weniger ins Gewicht, solange es nur um den Vergleich gleichzeitiger oder zeitlich nahe beieinanderliegender Intensitäten geht. Die Erfahrung hat jedoch gezeigt, dass sich Deklination und Inklination an einem gegebenen Ort fortwährend ändern und dass diese Änderungen nach vielen Jahren sehr groß werden. Es kann daher kein Zweifel daran bestehen, dass auch die Intensität des Erdmagnetismus ähnlichen, gleichsam säkularen Änderungen unterliegt. Gerade für \textit{diese} Frage ist das Verfahren offenkundig völlig wirkungslos. Dennoch wäre es für den Fortschritt der Naturwissenschaft höchst wünschenswert, gerade diese sehr wichtige Frage möglichst vollständig zu erhellen. Das kann gewiss nur gelingen, wenn das rein vergleichende Verfahren durch eines ersetzt wird, das von zufälligen Unterschieden der Nadeln völlig unabhängig ist und die Intensität des Erdmagnetismus auf feste Einheiten und absolute Maße zurückführt.

Die theoretischen Grundsätze, auf denen ein solches seit Langem gewünschtes Verfahren beruhen muss, lassen sich leicht aufstellen. Die Anzahl der Schwingungen, die die Nadel in einer gegebenen Zeit ausführt, hängt sowohl von der Intensität des Erdmagnetismus als auch von der Beschaffenheit der Nadel ab, nämlich vom statischen Moment der in ihr enthaltenen Elemente des freien Magnetismus und von ihrem Trägheitsmoment. Da sich dieses Trägheitsmoment \SourceLeaf{87}ohne Schwierigkeit bestimmen lässt, liefert uns die Beobachtung der Schwingungen offenbar das Produkt aus der Intensität des Erdmagnetismus und dem statischen Moment des Magnetismus der Nadel. Die beiden Größen lassen sich nur trennen, wenn wir Beobachtungen anderer Art heranziehen, in denen sie anders miteinander verknüpft sind. Dazu nehme man eine zweite Nadel hinzu und setze sie sowohl der Wirkung des Erdmagnetismus als auch der des Magnetismus der ersten Nadel aus, um das Verhältnis dieser beiden Wirkungen zu ermitteln. Beide Wirkungen hängen zwar von der Verteilung des freien Magnetismus in der zweiten Nadel ab. Die zweite hängt außerdem von der Beschaffenheit der ersten Nadel, vom Abstand der Mittelpunkte, von der Lage ihrer Verbindungslinie gegenüber den magnetischen Achsen beider Nadeln und schließlich vom Gesetz der magnetischen Anziehung und Abstoßung ab. Der unvergessliche \textsc{Tobias Mayer} hatte als Erster vermutet, dass dieses Gesetz insofern mit dem Gravitationsgesetz übereinstimmt, als auch diese Wirkungen umgekehrt proportional zum Quadrat der Entfernung abnehmen. Die Versuche der bedeutenden Physiker \textsc{Coulomb} und \textsc{Hansteen} haben dieser Vermutung große Wahrscheinlichkeit verliehen, und die neuesten Versuche erheben sie über jeden Zweifel. Man muss jedoch sorgfältig darauf achten, dass sich dieses Gesetz auf die einzelnen Elemente des freien Magnetismus bezieht. Die Gesamtwirkung eines magnetisierten Körpers verhält sich ganz anders. Bei sehr großen Entfernungen wird sie, wie sich aus eben diesem Gesetz herleiten lässt, dem Kehrwert der dritten Potenz der Entfernung annähernd proportional. Die mit der dritten Potenz der Entfernung multiplizierte Wirkung der Nadel nähert sich somit bei stetig wachsender Entfernung und sonst gleichen Umständen asymptotisch einem konstanten Wert. Drückt man die Entfernungen mit einer beliebig gewählten Länge als Einheit durch Zahlen aus, so hat dieser Wert dieselbe Dimension wie die Wirkung der Erdkraft und lässt sich mit ihr vergleichen. Der Grenzwert dieses Verhältnisses ist durch geeignete Anordnung und Durchführung der Versuche zu bestimmen. Da er nur das statische Moment des Magnetismus der ersten Nadel enthält, erhält man damit den Quotienten aus diesem Moment und der Intensität der Erdkraft. Vergleicht man ihn mit dem zuvor ermittelten Produkt dieser Größen, so lässt sich das statische Moment eliminieren, und man erhält den Wert der Intensität des Erdmagnetismus.

Um die Wirkungen des Erdmagnetismus und der ersten Nadel auf die zweite Nadel im Versuch zu untersuchen, stehen zwei Wege offen. Wir können die zweite Nadel entweder in Bewegung oder im Gleichgewicht beobachten. Beim ersten Verfahren werden ihre Schwingungen beobachtet, während sich die Wirkung des Erdmagnetismus mit der Wirkung der ersten Nadel verbindet. Die erste Nadel wird in geeignetem Abstand so angebracht, dass ihre Achse in dem magnetischen Meridian liegt, der durch den Mittelpunkt der schwingenden Nadel geht. Die Schwingungen werden so beschleunigt oder verzögert, je nachdem, ob befreundete oder feindliche Pole \SourceLeaf{88}einander zugekehrt sind. Man kann die Schwingungsdauern für die beiden Lagen der ersten Nadel miteinander vergleichen oder die Schwingungsdauer für eine der beiden Lagen mit derjenigen, die sich bei entfernter erster Nadel allein unter der Wirkung des Erdmagnetismus ergibt. Daraus ergibt sich das Verhältnis dieser Kraft zur Wirkung der ersten Nadel. Beim zweiten Verfahren stellt man die erste Nadel so auf, dass die Richtung der Kraft, die sie im Bereich der frei aufgehängten zweiten Nadel ausübt, mit dem magnetischen Meridian einen Winkel, etwa einen rechten, bildet. Die zweite Nadel wird dadurch aus dem magnetischen Meridian abgelenkt, und aus der Größe der Ablenkung lässt sich das Verhältnis der erdmagnetischen Kraft zur Wirkung der ersten Nadel erschließen.

Das erste Verfahren stimmt im Wesentlichen mit dem überein, das der bedeutende \textsc{Poisson} vor einigen Jahren vorgeschlagen hat. Doch die Versuche, die einige Physiker nach seiner Anleitung unternommen haben und die mir bekannt geworden sind, sind entweder ganz fehlgeschlagen oder haben nur eine grobe Annäherung ergeben.

Die Hauptschwierigkeit ergibt sich aus Folgendem. Aus den in mäßigen Entfernungen beobachteten Wirkungen der Nadel muss ein Grenzwert berechnet werden, der sich auf eine gleichsam unendlich große Entfernung bezieht. Je mehr Unbekannte zu eliminieren sind, die von der individuellen Beschaffenheit der Nadeln abhängen, desto mehr werden die dafür erforderlichen Eliminationen durch die kleinsten Beobachtungsfehler gestört, ja sogar unbrauchbar gemacht. Die Anzahl der Unbekannten lässt sich aber nur dann auf sehr wenige verringern, wenn die Wirkungen in Entfernungen beobachtet wurden, die im Verhältnis zur Länge der Nadeln hinreichend groß sind, sodass die Wirkungen selbst schon sehr klein geworden sind. Um so kleine Wirkungen genau zu messen, reichen die bisher verwendeten praktischen Hilfsmittel jedoch nicht aus.

Ich sah daher, dass ich vor allem neue Hilfsmittel entwickeln musste, mit denen ich sowohl die Schwingungsdauern als auch die Richtungen der Nadeln mit weit größerer Genauigkeit als bisher beobachten und messen konnte. Die dazu unternommenen und mehrere Monate fortgesetzten Arbeiten, bei denen mich der hervorragende \textsc{Weber} auf vielfältige Weise unterstützte, führten zum gewünschten Ziel, erfüllten die Erwartungen und übertrafen sie bei Weitem. Um die Genauigkeit der Versuche der Genauigkeit astronomischer Beobachtungen gleichzubringen, fehlt nur noch ein Ort, der vor dem Einfluss von Eisen in der Nähe und vor Luftbewegungen völlig geschützt ist. Es stehen zwei Geräte zur Verfügung, die sich ebenso durch ihre Einfachheit wie durch die Genauigkeit auszeichnen, die sie bieten. Ihre Beschreibung muss ich mir für eine andere Gelegenheit vorbehalten. In dieser Abhandlung lege ich den Physikern die Versuche vor, die wir zur Bestimmung der Intensität des Erdmagnetismus bisher in unserer Sternwarte durchgeführt haben.
